% registry: BOLCHOUN
\methodhead{BOLCHOUN}{Bolchoun, Baumgartner and Kaufmann: substitute amplitude and out-of-phase damage}\label{sheet:BOLCHOUN}

\block{Source} A.~Bolchoun, J.~Baumgartner, H.~Kaufmann, \emph{Int.\ J.\ Fatigue} (2017),
accepted manuscript \cite{bolchoun2017}, Eqs.~(4)--(15), Sections~2--5 and~8.

\block{Idea and scope} Integral approach with its own life law. A cycle is expressed by one
substitute amplitude, the integral of the normal stress amplitudes over all planes. Its in-phase
damage is read on a W\"ohler curve interpolated between the axial and torsion curves by the ratio
of the largest shear and normal amplitudes. The effect of rotating shear vectors is measured by
the area they sweep over all planes, $F_{OP}$, and turned into an extra damage calibrated on one
experimental out-of-phase curve. Proposed for variable-amplitude in-phase and out-of-phase
loading of thin-walled welded magnesium tubes (notch stress approach, $r_{ref}=0.05$\,mm).

\block{Inputs}
\begin{center}\small
\begin{tabular}{@{}llp{100mm}@{}}
\toprule
Symbol & Unit & Meaning\\\midrule
$\mathbf n$ & -- & plane normal on the unit hemisphere, $d\mathbf n=\cos\psi\,d\theta\,d\psi$\\
$\sigma_{n,a}$ & MPa & normal stress amplitude on the plane, $\sqrt{2\,\mathrm{Var}(\sigma_n)}$\\
$\boldsymbol\tau_{n,0}(t)$ & MPa & shear stress vector on the plane minus its mean\\
$\tau_{max,a}$, $\sigma_{max,a}$ & MPa & largest shear and normal stress amplitudes\\
$\sigma_f,b$; $\tau_f,b_0$ & MPa, -- & Basquin constants of the R $=-1$ axial and torsion curves\\
$\sigma_{f,OP},b_{OP}$ & MPa, -- & Basquin constants of the out-of-phase curve (axial amplitude)\\
$\lambda,\delta$ & --, deg & ratio $\tau_a/\sigma_a$ and phase of the out-of-phase path\\
$\sigma_{a,s}$ & MPa & substitute stress amplitude, Eq.~\eqref{eq:bol-sas}\\
\bottomrule
\end{tabular}
\end{center}

\block{Formulation} Out-of-phase measure (Eq.~9), interpolation factor and curve (Eq.~11) and
substitute amplitude (Eq.~12):
\begin{gather}
F_{OP}=\frac12\int_{t_1}^{t_2}\!\!\iint_{\mathbf n}\bigl\|\boldsymbol\tau_{n,0}\times\dot{\boldsymbol\tau}_{n,0}\bigr\|\,d\mathbf n\,dt,\qquad
f_I=\frac{\tau_{max,a}}{\sigma_{max,a}},\quad\tfrac12\le f_I\le1,\\
X=2\Bigl[X_{ax}(1-f_I)+X_{tors}\bigl(f_I-\tfrac12\bigr)\Bigr],\quad X=k,\ N_k,\ L_k,\qquad
\sigma_{a,s}=\iint_{\mathbf n}\sigma_{n,a}\,d\mathbf n .\label{eq:bol-sas}
\end{gather}
The out-of-phase damage (Eq.~14) compares the experimental out-of-phase curve with the in-phase
curve interpolated at its $f_{I,OP}$; the damage of a cycle is Eq.~(15), or its minus-sign variant
(Section~8) when out-of-phase loading is less damaging:
\begin{equation}
D_{OP}\bigl(F_{OP}(\sigma_{a,s})\bigr)=\sqrt{D_{total}^2-D_{IP}^2},\qquad
D_{total,C}=\sqrt{D_{IP,C}^2\pm\bigl(D_{OP,C}^{\,k_{I,C}/k_{I,OP}}\bigr)^2}.
\end{equation}
For one harmonic, Eq.~(9) is the area of the shear ellipses,
$F_{OP}=\iint_{\mathbf n}\pi\,\|\boldsymbol\tau_s\times\boldsymbol\tau_c\|\,d\mathbf n$.

\block{Constants and calibration} Axial, torsional and one out-of-phase W\"ohler curve in
$\sigma_{a,s}$. In the repository the axial and torsion curves are the Basquin fits on the R $=-1$
failures (workbook values first), converted by the unit factors $2\pi/3$ (tension) and $8/3$
(torsion). The out-of-phase curve is the Basquin fit on the most populated fully reversed
out-of-phase tension--torsion path of the dataset (5 tests, 1 decade, slope distinct from 0); its
tests are design points and the model is skipped without such a path.

\block{Life} Native. Paper: Palmgren--Miner with Haibach's slope $2k-2$ below the knee,
$N=L_S D_{al}/D_{total}$. Repository: constant amplitude, $N=1/D_{total,C}$.

\block{Repository choices}
\begin{itemize}
\item Constant amplitude, one harmonic: no rainflow counting, no Miner sum, no Haibach slope,
  $D_{al}=1$.
\item Basquin curves without a knee; Eq.~(11) interpolates $k=-1/b$ and $L_k$ at $N_k=10^6$
  (maintainer decision).
\item Out-of-phase curve fitted on the dataset; $D_{OP}(F)$ evaluated exactly on that path at the
  amplitude with the same $F_{OP}$ (maintainer decision).
\item Minus-sign variant when $D_{total}<D_{IP}$ on the out-of-phase curve, decided at the
  amplitude of each cycle; $D_{total,C}^2\le0$ is an infinite life (maintainer decision). On the
  calibration path Eq.~(15) returns the out-of-phase curve exactly; far above the in-phase curve
  the two terms nearly cancel and a small change of $f_I$ can give an infinite life.
\item Mean stresses out of scope, as in the paper (maintainer decision).
\item $\tau_{max,a}$ is the major semi-axis of the shear ellipse (the Maximum Variance Method
  value for one harmonic); maxima refined on a $48\times96$ sphere grid, also used for the
  hemisphere integrals. $f_I<\tfrac12$ uses the axial curve, $f_I>1$ the torsion curve.
\end{itemize}

\block{Worked example} Illustrative steel: $\sigma_f=1154$\MPa, $b=-0.061$, $\tau_f=940$\MPa,
$b_0=-0.065$; out-of-phase curve $\sigma_{f,OP}=773$\MPa, $b_{OP}=-0.036$ on a $90^\circ$ path with
$\lambda=0.577$. At $\sigma_a=500$\MPa: tension $\sigma_{a,s}=1047.2$\MPa, $f_I=0.5$,
$N=450\,459$; in phase with $\tau_a=0.577\sigma_a$, $\sigma_{a,s}=1244.5$\MPa, $f_I=0.604$,
$N=25\,784$; at $90^\circ$, $\sigma_{a,s}=1403.2$\MPa, $F_{OP}=356\,172$\,MPa$^2$, $f_I=0.577$,
$N=90\,113$, the out-of-phase curve itself (minus-sign variant).

\block{Verification} Checked: the unit factors of Eq.~(12); Eq.~(9) against the time integral and
zero for proportional loading; $f_I$ end points and its proportional closed form; Eq.~(11) end
points; tension and torsion return their Basquin curves; Eqs.~(14)--(15) return the out-of-phase
curve on its own path with either sign. The paper tabulates no per-test value for a reproducible
stress state.
